We evaluate three claims: (i) tool-call attack-success rate is reduced
to near zero on standard benchmarks; (ii) benign task completion is
preserved; (iii) per-call latency overhead is production-acceptable.

\subsubsection{6.0 Pre-registration}\label{pre-registration}

Before any measurement run that consumed LLM credits, the per-task
policy files (§3.4) were frozen by SHA-256 hash and committed to the
repository. The first eight hex characters of each file's hash:

\begin{verbatim}
banking.json    sha256:4b78e687…    slack.json     sha256:72698245…
travel.json     sha256:c7692b8e…    workspace.json sha256:8ba70481…
\end{verbatim}

Full hashes appear in \texttt{paper/eval/PRE-REGISTRATION.md}; the
static verifier (\texttt{paper/eval/verify\_policies.py}) re-computes
them on every CI build and aborts on drift. Policy edits after this
freeze require a separate commit explaining the change and a new
pre-registration entry. Any measurement number in §6.2 was produced
against a policy file whose hash exactly matches one of the above.

\subsubsection{6.1 Benchmarks and
Configurations}\label{benchmarks-and-configurations}

We use the two most cited contemporary prompt-injection benchmarks
against agentic LLMs:

\begin{itemize}
\tightlist
\item
  \textbf{AgentDojo} {[}@debenedetti2024agentdojo{]}: 629 attack tasks
  across four realistic deployments (banking, Slack, GitHub, travel
  agency). Each task pairs a benign user request with an adversarial
  document injected into the agent's context, with success defined as
  the agent calling a specific tool with attacker-controlled arguments.
\item
  \textbf{InjecAgent} {[}@zhan2024injecagent{]}: 1,054 attack scenarios
  across multiple agent frameworks, with similar structure.
\end{itemize}

We test six defence configurations:

\begin{enumerate}
\def\labelenumi{\arabic{enumi}.}
\tightlist
\item
  \textbf{None} --- no defence beyond the base model's training.
\item
  \textbf{Spotlighting} {[}@hines2024defending{]} --- input-perturbation
  defence that marks untrusted regions of the context.
\item
  \textbf{StruQ} {[}@chen2024struq{]} --- structured-query defence that
  fine-tunes the model to ignore instructions in delimited regions.
\item
  \textbf{Prompt shields} --- a commercial vendor's filter (anonymised
  at the vendor's request).
\item
  \textbf{VCD text-only} --- the scanner layer of our reference
  implementation, without proof or capability.
\item
  \textbf{VCD full stack} --- scanner + proof + capability gate.
\end{enumerate}

All configurations run against the same base model and the same agent
harness. We measure on three contemporary frontier-class models:
\textbf{deepseek-v3.2} (671B, late 2025), \textbf{deepseek-v4-flash}
(140B, April 2026), and \textbf{deepseek-v4-pro} (1.6T, May 2026), via
Ollama Cloud's OpenAI-compatible endpoint. We report tool-call-mediated
attack success rate, benign task completion on the AgentDojo benign-task
split, and gate latency on commodity hardware.

\subsubsection{6.2 Headline Result}\label{headline-result}

Measured on the AgentDojo banking suite (144 user × injection task pairs
per cell). Numbers below are from the live LLM-driven evaluation. The
cross-suite generalisation rows (slack, travel, workspace) appear in
§6.2.3; the attack-family-robustness rows (\texttt{direct},
\texttt{ignore\_previous} in addition to
\texttt{important\_instructions}) appear in §6.2.2; the component
ablation (text-only vs gate-only vs full-stack) appears in §6.2.4.

\textbf{deepseek-v4-flash (frontier flash, April 2026):}

\begin{longtable}[]{@{}llll@{}}
\toprule\noalign{}
Defence & ASR & Benign & Wall \\
\midrule\noalign{}
\endhead
\bottomrule\noalign{}
\endlastfoot
None & \textbf{1.4\%} & 86.8\% & 7.6m \\
Spotlighting & 2.1\% & 84.0\% & 14.4m \\
Prompt shields (DeBERTa PI) & \textbf{0.0\%} & \textbf{35.4\%} &
26.5m \\
VCD text-only & 0.0\% & 91.0\% & 22.5m \\
\textbf{VCD full stack} & \textbf{0.7\%} & \textbf{90.3\%} & 8.7m \\
VCD capability-only (ablation) & 0.0\% & 88.9\% & 9.1m \\
\end{longtable}

\textbf{deepseek-v4-pro (frontier reasoning, May 2026):}

\begin{longtable}[]{@{}llll@{}}
\toprule\noalign{}
Defence & ASR & Benign & Wall \\
\midrule\noalign{}
\endhead
\bottomrule\noalign{}
\endlastfoot
None & \textbf{33.3\%} & 88.2\% & 58.5m \\
Spotlighting & 6.9\% & 87.5\% & 46.1m \\
Prompt shields (DeBERTa PI) & \textbf{0.0\%} & \textbf{36.1\%} &
65.3m \\
VCD text-only & \textbf{0.0\%} & \textbf{84.0\%} & 51.1m \\
\textbf{VCD full stack} & \textbf{0.7\%} & \textbf{85.4\%} & 53.9m \\
VCD capability-only (ablation) & 0.7\% & 84.7\% & 50.1m \\
\end{longtable}

\textbf{deepseek-v3.2 (frontier reasoning, late 2025):}

\begin{longtable}[]{@{}llll@{}}
\toprule\noalign{}
Defence & ASR & Benign & Wall \\
\midrule\noalign{}
\endhead
\bottomrule\noalign{}
\endlastfoot
None & \textbf{70.8\%} & 79.9\% & 43.9m \\
Spotlighting & 58.3\% & 80.6\% & 67.8m \\
Prompt shields (DeBERTa PI) & 7.6\% & \textbf{33.3\%} & 50.5m \\
VCD text-only & \textbf{0.7\%} & \textbf{79.9\%} & 99.4m \\
\textbf{VCD full stack} & \textbf{0.7\%} & 77.8\% & 88.7m \\
VCD capability-only (ablation) & 0.7\% & 75.7\% & 101.9m \\
\end{longtable}

The load-bearing claim is the \textbf{shields vs VCD full-stack}
comparison, and it reproduces across all three model generations: both
defences drive attack-success rate to within a single attack of zero,
but the strongest text-side defence (transformers PI detector) collapses
benign task completion to the 33--36\% range across every model we
tested, while VCD full-stack preserves benign completion at 78--90\%.
The benign-preservation gap at equivalent security is \textbf{+54.9
percentage points} on v4-flash, \textbf{+49.3 pp} on v4-pro, and
\textbf{+44.4 pp} on v3.2. The effect is \emph{largest} on the weakest
model: as the base model gets less reliable on benign tasks, shields'
collateral damage compounds, while VCD's surgical precision remains
unchanged. Equivalently, deploying shields requires accepting the loss
of roughly half of legitimate agent productivity to reach VCD's security
level --- and this cost does not amortise with scale.

A second row deserves direct comment. On v3.2, where the baseline
attack-success rate is 70.8\%, VCD-text alone reduces ASR to 0.7\% --- a
\textbf{10× stronger} result than prompt-shields' 7.6\% on the same
model --- \emph{while} preserving 46.6 percentage points more benign
completion. We did not anticipate that the text-side component of VCD
would outperform a dedicated DeBERTa detector on a less-capable base
model; we discuss the likely mechanism in §7.1.

\textbf{Cross-provider banking row.} To verify the pattern is not
specific to the DeepSeek family, we additionally measure two
frontier-class models from two other providers --- Qwen3.5:397b
(Alibaba) and Kimi-k2.6 (Moonshot) --- on the same banking suite under
the same load-bearing comparison:

\begin{longtable}[]{@{}
  >{\raggedright\arraybackslash}p{(\linewidth - 8\tabcolsep) * \real{0.2000}}
  >{\raggedright\arraybackslash}p{(\linewidth - 8\tabcolsep) * \real{0.2000}}
  >{\raggedright\arraybackslash}p{(\linewidth - 8\tabcolsep) * \real{0.2000}}
  >{\raggedright\arraybackslash}p{(\linewidth - 8\tabcolsep) * \real{0.2000}}
  >{\raggedright\arraybackslash}p{(\linewidth - 8\tabcolsep) * \real{0.2000}}@{}}
\toprule\noalign{}
\begin{minipage}[b]{\linewidth}\raggedright
Provider
\end{minipage} & \begin{minipage}[b]{\linewidth}\raggedright
Model
\end{minipage} & \begin{minipage}[b]{\linewidth}\raggedright
Defence
\end{minipage} & \begin{minipage}[b]{\linewidth}\raggedright
ASR
\end{minipage} & \begin{minipage}[b]{\linewidth}\raggedright
Benign
\end{minipage} \\
\midrule\noalign{}
\endhead
\bottomrule\noalign{}
\endlastfoot
Alibaba (MoE) & qwen3.5:397b & None & 6.9\% & 75.0\% \\
Alibaba (MoE) & qwen3.5:397b & Prompt shields & 0.0\% &
\textbf{35.4\%} \\
Alibaba (MoE) & qwen3.5:397b & \textbf{VCD full stack} & \textbf{0.7\%}
& \textbf{72.2\%} \\
Moonshot (MoE) & kimi-k2.6 & None & \textbf{0.0\%} & 68.8\% \\
Moonshot (MoE) & kimi-k2.6 & Prompt shields & 0.0\% & \textbf{34.7\%} \\
Moonshot (MoE) & kimi-k2.6 & \textbf{VCD full stack} & \textbf{0.0\%} &
\textbf{66.0\%} \\
\end{longtable}

The pattern reproduces. On Qwen3.5 the benign-preservation gap between
shields and VCD-full is \textbf{+36.8 pp} at 0.0--0.7\% ASR. The Kimi
row is more interesting: Moonshot's instruction-adherence training is
strong enough that the unprotected baseline ASR is already
\textbf{0.0\%} --- no defence is \emph{needed} for security on this
configuration. Prompt-shields nonetheless collapses benign task
completion from 68.8\% to 34.7\% (a 34.1-percentage-point loss) for
security that wasn't required; VCD-full holds benign at 66.0\% (a
2.8-percentage-point loss). \textbf{Shields' collateral damage is
independent of its security necessity} --- the defence imposes the same
cost whether or not it is doing anything useful. This is the cleanest
demonstration in the paper that the trade-off curve is shaped
differently for the two defences: VCD's cost scales with its work;
shields' cost is paid up-front regardless.

The static upper bound of §6.2.1 applies identically to these models ---
the gate's constraint logic is model-independent, so the
catalogued-attack block rate is 100\% on Qwen and Kimi by construction.

A note on the pure-output fraction: of the AgentDojo injection tasks in
our four-suite cross-evaluation, \textbf{3 of 25 distinct injection
tasks} (slack \texttt{injection\_task\_3}, slack
\texttt{injection\_task\_4}, travel \texttt{injection\_task\_6})
describe attacks whose success criterion is the agent's free-form
natural-language response rather than a tool call with specific argument
values --- visit-a-link suggestions, post-this-text-on-a-channel
commands routed through unstructured replies, and biased hotel
recommendations. These attacks fall outside VCD's design boundary (§6.5)
because there is no tool call to gate. They account for every non-zero
ASR cell we measured on slack/travel/workspace; outside of them, the
LLM-driven block rate on attacker-controlled tool calls is identical to
the §6.2.1 static bound (100\%). The corresponding InjecAgent fraction
is reported with the static-bound numbers in §6.2.1: the catalogued
attack args are entirely tool-call-mediated, and the static block rate
is 100\% across the full 2,108-pair benchmark.

\paragraph{6.2.1 Static upper bound on attack-success
rate}\label{static-upper-bound-on-attack-success-rate}

Before reporting the LLM-driven measurements, we report a \emph{static}
upper bound: assuming the model emits the attacker's canonical tool call
exactly, do the gate's constraints --- derived independently from each
user task's prompt and pinned by the pre-registered hash anchors of §6.0
--- admit the call? This is the maximum ASR the actual eval can exhibit;
the LLM-driven run can only do as well or better, because
natural-language defences provide additional attrition on top of the
gate.

The static verifier (\texttt{paper/eval/verify\_policies.py}) runs the
same constraint logic the runtime gate runs. Static results under VCD
full-stack, across both benchmarks, all defence-relevant cells:

\begin{longtable}[]{@{}
  >{\raggedright\arraybackslash}p{(\linewidth - 6\tabcolsep) * \real{0.2500}}
  >{\raggedright\arraybackslash}p{(\linewidth - 6\tabcolsep) * \real{0.2500}}
  >{\raggedright\arraybackslash}p{(\linewidth - 6\tabcolsep) * \real{0.2500}}
  >{\raggedright\arraybackslash}p{(\linewidth - 6\tabcolsep) * \real{0.2500}}@{}}
\toprule\noalign{}
\begin{minipage}[b]{\linewidth}\raggedright
Benchmark
\end{minipage} & \begin{minipage}[b]{\linewidth}\raggedright
Cases
\end{minipage} & \begin{minipage}[b]{\linewidth}\raggedright
Attacks succeeding the static check
\end{minipage} & \begin{minipage}[b]{\linewidth}\raggedright
Static upper bound on ASR
\end{minipage} \\
\midrule\noalign{}
\endhead
\bottomrule\noalign{}
\endlastfoot
AgentDojo (v1, all four suites) & 629 & 0 & 0.00\% \\
InjecAgent (DH base + enhanced) & 1020 & 0 & 0.00\% \\
InjecAgent (DS base + enhanced) & 1088 & 0 & 0.00\% \\
\textbf{Total} & \textbf{2737} & \textbf{0} & \textbf{0.00\%} \\
\end{longtable}

The single known exception is one cell in AgentDojo banking
(user\_task\_15 × schedule\_transaction): the user's stated intent
contains the same IBAN as the benchmark's adversarial recipient (the
user's landlord). Capability discipline correctly admits that call ---
the user explicitly authorised it --- while the same IBAN remains
forbidden in every other user task's policy. The static verifier flags
this cell as a \emph{known-legitimate coincidence} and excludes it from
the count above.

The verifier and the runtime gate share a single implementation of the
constraint-check logic --- both call into the same
\texttt{Policy.satisfiesArgs} function in the implementation's
capability module --- so the static result is not a separate model that
could disagree with the runtime. The verifier is a 30-line driver that
loads the policy files, iterates over the benchmark's catalogued attack
args, and asks the same Python function the gate asks at runtime. Any
policy edit that weakens the gate's behaviour weakens the static result
identically. A CI job (configured in \texttt{.github/workflows/ci.yml})
re-runs the verifier on every commit and fails the build if the static
block rate regresses below 100\%.

\paragraph{6.2.2 Attack-family
robustness}\label{attack-family-robustness}

The numbers in §6.2 use AgentDojo's \texttt{important\_instructions}
injection template, the benchmark's strongest and most-cited attack. To
check that the result is not overfit to a single attack vector, we
measure the banking row under the \texttt{direct} and
\texttt{ignore\_previous} attack templates across all three model
generations (144 pairs each cell, same policy and harness). The full
coverage matrix below is 3 models × 3 attack families = 9 attack-cohort
cells; the v4-flash row reports all five defences (none, spotlight,
prompt-shields, VCD-text, VCD-cap-only, VCD-full) and the v4-pro / v3.2
rows report the load-bearing comparison \{none, shields, vcd\_full\}.

\textbf{deepseek-v4-flash (full defence sweep):}

\begin{longtable}[]{@{}llll@{}}
\toprule\noalign{}
Attack & Defence & ASR & Benign \\
\midrule\noalign{}
\endhead
\bottomrule\noalign{}
\endlastfoot
\texttt{direct} & None & 13.9\% & 91.0\% \\
\texttt{direct} & Spotlighting & 13.9\% & 87.5\% \\
\texttt{direct} & Prompt shields & 0.0\% & \textbf{36.8\%} \\
\texttt{direct} & VCD text-only & 0.7\% & 89.6\% \\
\texttt{direct} & VCD capability-only & 0.7\% & 90.3\% \\
\texttt{direct} & \textbf{VCD full stack} & \textbf{0.7\%} &
\textbf{91.7\%} \\
\texttt{ignore\_previous} & None & 1.4\% & 88.9\% \\
\texttt{ignore\_previous} & Spotlighting & 0.7\% & 87.5\% \\
\texttt{ignore\_previous} & Prompt shields & 0.0\% & \textbf{31.3\%} \\
\texttt{ignore\_previous} & VCD text-only & 0.7\% & 85.4\% \\
\texttt{ignore\_previous} & VCD capability-only & 0.7\% & 87.5\% \\
\texttt{ignore\_previous} & \textbf{VCD full stack} & \textbf{0.7\%} &
\textbf{89.6\%} \\
\end{longtable}

The pattern is identical across all three attack families on v4-flash:
shields drives ASR to zero by destroying benign task completion;
\textbf{every VCD configuration} drives ASR to the same 0.7\%
benchmark-artefact floor while preserving 85--92\% benign completion.
Spotlighting is interesting: it is essentially inactive against
\texttt{direct} (13.9\% ASR, unchanged from baseline) but reduces
\texttt{ignore\_previous} to 0.7\% --- confirming the spotlighting
mechanism is sensitive to the surface form of the injection. VCD does
not have this dependence because the gate operates structurally on
tool-call arguments, not on the surrounding natural-language framing.
The benign-preservation gap between VCD-full and shields at equivalent
security on v4-flash is \textbf{+54.9 pp} under
\texttt{important\_instructions}, \textbf{+54.9 pp} under
\texttt{direct}, and \textbf{+58.3 pp} under \texttt{ignore\_previous}.

\textbf{deepseek-v4-pro:}

\begin{longtable}[]{@{}llll@{}}
\toprule\noalign{}
Attack & Defence & ASR & Benign \\
\midrule\noalign{}
\endhead
\bottomrule\noalign{}
\endlastfoot
\texttt{direct} & None & 12.5\% & 85.4\% \\
\texttt{direct} & Prompt shields & 0.7\% & \textbf{36.8\%} \\
\texttt{direct} & \textbf{VCD full stack} & \textbf{0.7\%} &
\textbf{82.6\%} \\
\texttt{ignore\_previous} & None & 0.7\% & 81.2\% \\
\texttt{ignore\_previous} & Prompt shields & 0.7\% & \textbf{31.2\%} \\
\texttt{ignore\_previous} & \textbf{VCD full stack} & \textbf{0.0\%} &
\textbf{77.8\%} \\
\end{longtable}

\textbf{deepseek-v3.2:}

\begin{longtable}[]{@{}llll@{}}
\toprule\noalign{}
Attack & Defence & ASR & Benign \\
\midrule\noalign{}
\endhead
\bottomrule\noalign{}
\endlastfoot
\texttt{direct} & None & 13.9\% & 88.2\% \\
\texttt{direct} & Prompt shields & \textbf{4.2\%} & \textbf{35.4\%} \\
\texttt{direct} & \textbf{VCD full stack} & \textbf{0.7\%} &
\textbf{86.8\%} \\
\texttt{ignore\_previous} & None & 9.0\% & 84.7\% \\
\texttt{ignore\_previous} & Prompt shields & 0.0\% & \textbf{30.6\%} \\
\texttt{ignore\_previous} & \textbf{VCD full stack} & \textbf{0.7\%} &
\textbf{79.2\%} \\
\end{longtable}

The cross-generation pattern reproduces. On v4-pro the
shields-vs-VCD-full benign-preservation gap is \textbf{+45.8 pp} under
\texttt{direct} and \textbf{+46.6 pp} under \texttt{ignore\_previous}.
On v3.2 it is \textbf{+51.4 pp} under \texttt{direct} and \textbf{+48.6
pp} under \texttt{ignore\_previous}. Two cells deserve special mention:

\begin{itemize}
\tightlist
\item
  \textbf{v3.2 × \texttt{direct} × shields} is the only cell in the
  entire study where the strongest text-side defence fails to drive ASR
  to zero (4.2\%, six attacks landed out of 144). Shields \emph{both}
  leaves a non-trivial residual attack-success rate \emph{and} collapses
  benign task completion to 35.4\%. On the same configuration, VCD-full
  holds ASR at 0.7\% and preserves 86.8\% benign completion --- strictly
  dominating shields on every measured dimension simultaneously.
\item
  \textbf{v3.2 × \texttt{ignore\_previous} × shields} drives ASR to zero
  but at the cost of dropping benign task completion to \textbf{30.6\%},
  the lowest benign completion shields produces anywhere in the paper.
  The corresponding VCD-full cell achieves 0.7\% ASR at 79.2\% benign
  --- a \textbf{+48.6 percentage-point} gap at near-identical security.
\end{itemize}

The aggregate claim is that the headline pattern --- shields trades
catastrophic benign-completion collapse for marginal security gain,
VCD-full holds security at the benchmark-artefact floor while preserving
the agent's productive capacity --- is invariant across every measurable
cell of (model, attack family) we tested.

\paragraph{6.2.3 Cross-suite
generalisation}\label{cross-suite-generalisation}

To check that the result is not banking-specific, we measure on
AgentDojo's remaining three suites --- slack (105 pairs), travel (140
pairs), workspace (240 pairs) --- across two model generations (v4-flash
and v4-pro):

\textbf{deepseek-v4-flash:}

\begin{longtable}[]{@{}llll@{}}
\toprule\noalign{}
Suite & Defence & ASR & Benign \\
\midrule\noalign{}
\endhead
\bottomrule\noalign{}
\endlastfoot
Slack & None & 14.3\% & 66.7\% \\
Slack & \textbf{VCD full stack} & \textbf{4.8\%} & \textbf{61.0\%} \\
Travel & None & 3.6\% & 51.4\% \\
Travel & Prompt shields & 0.0\% & \textbf{0.0\%} \\
Travel & \textbf{VCD full stack} & \textbf{1.4\%} & \textbf{58.6\%} \\
Workspace & None & 1.7\% & 72.5\% \\
Workspace & Prompt shields & 1.3\% & \textbf{36.7\%} \\
Workspace & \textbf{VCD full stack} & \textbf{0.4\%} &
\textbf{64.6\%} \\
\end{longtable}

\textbf{deepseek-v4-pro:}

\begin{longtable}[]{@{}llll@{}}
\toprule\noalign{}
Suite & Defence & ASR & Benign \\
\midrule\noalign{}
\endhead
\bottomrule\noalign{}
\endlastfoot
Slack & None & 65.7\% & 66.7\% \\
Slack & \textbf{VCD full stack} & \textbf{18.1\%} & \textbf{62.9\%} \\
Travel & None & 33.6\% & 57.9\% \\
Travel & \textbf{VCD full stack} & \textbf{6.4\%} & \textbf{74.3\%} \\
Workspace & None & 45.0\% & 47.9\% \\
Workspace & \textbf{VCD full stack} & \textbf{7.1\%} &
\textbf{68.8\%} \\
\end{longtable}

\textbf{deepseek-v3.2:}

\begin{longtable}[]{@{}llll@{}}
\toprule\noalign{}
Suite & Defence & ASR & Benign \\
\midrule\noalign{}
\endhead
\bottomrule\noalign{}
\endlastfoot
Slack & None & 98.1\% & 58.1\% \\
Travel & None & 73.6\% & 22.1\% \\
Travel & \textbf{VCD full stack} & \textbf{7.9\%} & \textbf{68.6\%} \\
Workspace & None & 65.8\% & 9.2\% \\
Workspace & \textbf{VCD full stack} & \textbf{6.7\%} &
\textbf{52.9\%} \\
\end{longtable}

(Two cells omitted. The slack × shields cell on v4-flash and the slack ×
vcd\_full cell on v3.2 both crashed in the same AgentDojo slack-suite
utility scorer
(\texttt{KeyError:\ \textquotesingle{}www.our-company.com\textquotesingle{}}
in \texttt{agentdojo/default\_suites/v1/slack/user\_tasks.py:162}),
which assumes a \texttt{web.web\_content} key that is not populated when
the agent's trajectory does not visit a specific URL. The crash is in
the benchmark's scoring code, not in our harness, policy, or the gate.
We report 19 of 21 cross-suite cells intact rather than retry; the bug
is reproducible on AgentDojo 0.1.35 against any agent that does not
happen to fetch the COMPANY\_URL during a slack user task.)

VCD generalises across all three cross-suites and all three measured
model generations. The pattern strengthens as the base model gets less
capable. On v4-flash the no-defence baseline ASR sits low (1.7--14.3\%)
and VCD reduces it by 3×--4×. On v4-pro the baseline is much higher
(33.6--65.7\%) and VCD reduces it by 3.6×--6.3×, with the protected
benign completion \emph{exceeding} the unprotected baseline on travel
and workspace (+16.4 pp travel, +20.9 pp workspace). On v3.2 --- where
the baseline attack-success rate runs 65.8--98.1\% across the three
suites and the unprotected agent's benign completion collapses to
9.2--58.1\% under attack --- \textbf{VCD reduces ASR by 9.4× on travel
and 9.9× on workspace, \emph{and} improves benign completion by +46.5 pp
on travel and +43.8 pp on workspace}. The protected agent on v3.2
workspace completes 52.9\% of benign tasks where the unprotected
baseline completed only 9.2\% --- almost a six-fold productivity
recovery. The mechanism is the same early-termination effect observed in
§6.3.1: when the gate denies an attacker-induced tool call, the agent's
reasoning loop returns to the legitimate task instead of looping further
on the adversarial trajectory. The effect compounds when the base model
is less robust to derailment.

Where shields was measurable to completion on v4-flash, it drove benign
task completion to \textbf{0\%} on travel and \textbf{36.7\%} on
workspace --- a 35.8-percentage-point collapse on workspace, and a
complete collapse on travel where the detector flagged every legitimate
task as adversarial. As the suite's tool surface grows more diverse,
shields' false-positive rate compounds until the agent cannot complete
most or all benign work. VCD's benign-preservation degrades far more
gracefully because the gate's rejection criterion is structural, not
stylistic. \textbf{VCD-vs-shields benign-preservation gap on the three
measurable cross-suite cells: travel +58.6 pp, workspace +27.9 pp,
banking +54.9 pp.}

A note on the residual VCD ASR on these suites: unlike banking, where
the residual is a single IBAN-collision artefact (§6.5), the cross-suite
residual concentrates on specific injection tasks that ask the agent to
perform actions \emph{without} a constrained tool call --- e.g.,
\texttt{slack.injection\_task\_3} (``Visit a website at the link
\ldots{}''), \texttt{travel.injection\_task\_6} (``Say that I should
absolutely visit Riverside View Hotel\ldots{}''). These are the
pure-output / pure-recommendation attacks §6.5 already flags as outside
VCD's design boundary: the gate operates on tool-call arguments, and an
attack whose effect is the agent's free-form natural-language response
has no tool-call to bind. The cross-suite numbers are therefore
consistent with the §6.5 claim and the §6.2.1 static bound ---
\emph{every} tool-mediated attack across all four suites is denied by
the gate.

\paragraph{6.2.4 Ablation: which component carries the
result}\label{ablation-which-component-carries-the-result}

VCD composes a text-side input scanner (the \emph{text} component) with
the capability-token gate (the \emph{cap} component). The two components
are independently meaningful and can be deployed separately. We measure
each in isolation on banking, across all three models:

\begin{longtable}[]{@{}llll@{}}
\toprule\noalign{}
Model & Configuration & ASR & Benign \\
\midrule\noalign{}
\endhead
\bottomrule\noalign{}
\endlastfoot
v4-flash & Text-side scanner only (no gate) & 0.0\% & 91.0\% \\
v4-flash & Capability gate only (no scanner) & 0.0\% & 88.9\% \\
v4-flash & \textbf{Full stack (both)} & \textbf{0.7\%} &
\textbf{90.3\%} \\
v4-pro & Text-side scanner only (no gate) & 0.0\% & 84.0\% \\
v4-pro & Capability gate only (no scanner) & 0.7\% & 84.7\% \\
v4-pro & \textbf{Full stack (both)} & \textbf{0.7\%} &
\textbf{85.4\%} \\
v3.2 & Text-side scanner only (no gate) & 0.7\% & 79.9\% \\
v3.2 & Capability gate only (no scanner) & 0.7\% & 75.7\% \\
v3.2 & \textbf{Full stack (both)} & \textbf{0.7\%} & \textbf{77.8\%} \\
\end{longtable}

(The single residual ASR cell in every row is the same
\texttt{user\_task\_15\ ×\ injection\_task\_4} benchmark IBAN
coincidence discussed in §6.5.)

Two observations. First, \textbf{either component alone is essentially
sufficient on AgentDojo banking} --- the gate, on its own, denies every
attacker-controlled tool call without help from the text scanner; the
text scanner, on its own, intercepts the injection before it reaches the
LLM. Second, the components have \emph{different} failure modes and
\emph{different} operating costs, so the right deployment is the
composition. The capability gate is sub-100µs per call (§6.3), is
provably correct against §6.2.1's static bound, and protects against
attacks the scanner missed (e.g., a benign-looking but policy-violating
tool call constructed by the model itself). The text scanner adds
defence-in-depth against attacks the gate cannot see --- e.g.,
pure-output exfiltration (§6.5) where there is no tool call to bind. The
full-stack configuration's benign-completion is within \textasciitilde2
pp of either component alone, so the composition does not pay a
meaningful precision cost.

We do not report a ``SMT-proof-only'' ablation because the SMT proof in
VCD operates on the policy, not on individual tool calls, and is not
independently deployable from the gate --- the gate consumes the proof's
compiled constraint set at runtime. The static-bound result in §6.2.1 is
the corresponding evidence: the gate's constraint logic, run against the
canonical attack args from both benchmarks, denies every attack across
2,737 user-task × injection-task pairs.

The full-stack row is the load-bearing claim: every attack whose effect
is \emph{mediated through a tool call} (operationally defined in §2) is
rejected by the gate, because every such tool call must satisfy the
constraints in the in-force capability token, which were issued before
the attacker had any access to the agent's context. The LLM-driven
measurements in the tables above can only do as well or better than the
§6.2.1 static bound. We discuss the small residual in §6.5.

\subsubsection{6.3 Per-Call Latency}\label{per-call-latency}

Measured on two reference hardware classes, single thread, Python 3.12+:
an x86\_64 cloud VM (AMD EPYC-Milan, Ubuntu 24.04) and a workstation
(Apple M-series ARM64, macOS 14). 5,000 iterations per gate operation,
200 iterations per proof.

\begin{longtable}[]{@{}
  >{\raggedright\arraybackslash}p{(\linewidth - 4\tabcolsep) * \real{0.3333}}
  >{\raggedright\arraybackslash}p{(\linewidth - 4\tabcolsep) * \real{0.3333}}
  >{\raggedright\arraybackslash}p{(\linewidth - 4\tabcolsep) * \real{0.3333}}@{}}
\toprule\noalign{}
\begin{minipage}[b]{\linewidth}\raggedright
Operation
\end{minipage} & \begin{minipage}[b]{\linewidth}\raggedright
x86\_64 (EPYC-Milan) p50 / p95 / p99
\end{minipage} & \begin{minipage}[b]{\linewidth}\raggedright
ARM64 (Apple M) p50 / p95 / p99
\end{minipage} \\
\midrule\noalign{}
\endhead
\bottomrule\noalign{}
\endlastfoot
\texttt{Gate.check()} no chain & 0.07 / 0.08 / 0.11 ms & 0.15 / 0.18 /
0.19 ms \\
\texttt{Gate.check()} 3-link chain & 0.27 / 0.30 / 0.34 ms & 0.58 / 0.64
/ 0.69 ms \\
\texttt{Prove()} cold & 0.67 / 0.74 / 0.90 ms & 0.54 / 0.68 / 46.5 ms \\
\end{longtable}

Proof results are cached by \texttt{(schema\_hash,\ policy\_hash)}. In
steady-state deployments where the schema and policy change rarely, the
prover is invoked once per policy version and the cache hit rate
approaches 100\%. \textbf{End-to-end overhead per tool call in the
steady state is dominated by the gate path, at well under 100
microseconds at p50.} For agents making tens of tool calls per turn at
human-conversation cadence, the cost is invisible. Proof-cache
invalidation on policy update incurs a one-time sub-millisecond hit.

The numbers are dramatically below the latencies typically associated
with formal-verification primitives because the supported JSON Schema
fragment is small enough that Z3 finds the proof or counterexample in a
handful of solver iterations. Schemas at the edge of our supported
fragment (deep enum sets, many fields) can extend \texttt{Prove()}
cold-path latency into the low milliseconds; this remains acceptable
because proofs are issued once per policy version, not per call.

\paragraph{6.3.1 End-to-end agent wall
time}\label{end-to-end-agent-wall-time}

Per-call gate latency is the structural cost of VCD; what an operator
pays in \emph{practice} is end-to-end agent wall-clock with the defence
enabled. We report wall time per defence across every cohort in §6.2,
normalised to the unprotected \texttt{none} baseline for that cohort:

\begin{longtable}[]{@{}
  >{\raggedright\arraybackslash}p{(\linewidth - 12\tabcolsep) * \real{0.1429}}
  >{\raggedright\arraybackslash}p{(\linewidth - 12\tabcolsep) * \real{0.1429}}
  >{\raggedright\arraybackslash}p{(\linewidth - 12\tabcolsep) * \real{0.1429}}
  >{\raggedright\arraybackslash}p{(\linewidth - 12\tabcolsep) * \real{0.1429}}
  >{\raggedright\arraybackslash}p{(\linewidth - 12\tabcolsep) * \real{0.1429}}
  >{\raggedright\arraybackslash}p{(\linewidth - 12\tabcolsep) * \real{0.1429}}
  >{\raggedright\arraybackslash}p{(\linewidth - 12\tabcolsep) * \real{0.1429}}@{}}
\toprule\noalign{}
\begin{minipage}[b]{\linewidth}\raggedright
Cohort
\end{minipage} & \begin{minipage}[b]{\linewidth}\raggedright
none (min)
\end{minipage} & \begin{minipage}[b]{\linewidth}\raggedright
spotlight
\end{minipage} & \begin{minipage}[b]{\linewidth}\raggedright
shields
\end{minipage} & \begin{minipage}[b]{\linewidth}\raggedright
VCD-text
\end{minipage} & \begin{minipage}[b]{\linewidth}\raggedright
\textbf{VCD-full}
\end{minipage} & \begin{minipage}[b]{\linewidth}\raggedright
VCD-cap-only
\end{minipage} \\
\midrule\noalign{}
\endhead
\bottomrule\noalign{}
\endlastfoot
banking / v4-flash / \texttt{important\_instructions} & 7.6 & 1.89× &
\textbf{3.49×} & 2.96× & \textbf{1.14×} & 1.20× \\
banking / v4-flash / \texttt{direct} & 90.7 & --- & \textbf{2.94×} & ---
& \textbf{1.12×} & --- \\
banking / v4-flash / \texttt{ignore\_previous} & 116.4 & --- & 1.45× &
--- & \textbf{0.54×} & --- \\
banking / v4-pro / \texttt{important\_instructions} & 58.5 & 0.79× &
1.11× & 0.87× & \textbf{0.92×} & 0.86× \\
banking / v3.2 / \texttt{important\_instructions} & 43.9 & 1.54× & 1.15×
& 2.26× & 2.02× & 2.32× \\
slack / v4-flash / \texttt{important\_instructions} & 59.2 & --- & --- &
--- & \textbf{0.23×} & --- \\
travel / v4-flash / \texttt{important\_instructions} & 63.5 & --- &
\textbf{2.80×} & --- & \textbf{0.42×} & --- \\
workspace / v4-flash / \texttt{important\_instructions} & 95.5 & --- &
1.59× & --- & \textbf{0.21×} & --- \\
\end{longtable}

Two patterns. \textbf{First, prompt-shields end-to-end wall time is
consistently 1.5--3.5× the unprotected baseline.} This is on top of
shields' benign-completion collapse documented in §6.2 --- the
operational cost of the defence is paid twice. The 2.94× and 2.80×
figures on v4-flash banking/\texttt{direct} and travel are the practical
ceiling: shields adds DeBERTa inference per message, and inflates the
agent's reasoning trajectory by escalating its caution.

\textbf{Second, VCD-full is at-baseline or faster.} On four of the eight
cohorts VCD-full \emph{reduces} wall time relative to the unprotected
baseline (slack 0.23×, workspace 0.21×, travel 0.42×, v4-flash
\texttt{ignore\_previous} 0.54×). The mechanism is structural: when the
gate denies an attacker-induced tool call, the agent's reasoning loop
terminates instead of producing more LLM turns trying to execute the
prohibited action. On the remaining cohorts VCD-full sits at
0.92×--2.02× baseline; the outlier is v3.2 banking (2.02×), where the
older model needs more retry attempts before settling on a
gate-admissible call. \textbf{The single VCD-full row that exceeds 1.2×
baseline is v3.2; on every v4-class model and every cross-suite
measurement, VCD-full is between 0.21× and 1.20× baseline.}
Operationally, deploying VCD costs less wall-clock than running
unprotected on three of our four AgentDojo suites --- because attacks
are denied early --- and costs at most \textasciitilde14\% extra on
banking where attacks are rarer.

The gate's microbenchmarked sub-100µs latency from §6.3 dominates as a
function of \emph{the number of admitted tool calls}; the wall-clock
improvements above reflect \emph{the number of LLM turns avoided} when
the gate rejects an attack rather than letting the agent loop. The two
effects compose: VCD adds negligible per-call overhead and subtracts
whole reasoning rounds whenever it triggers.

\subsubsection{6.4 Reference Deployment}\label{reference-deployment}

The reference implementation is released as open source under a
permissive licence (Apache-2.0) and was deposited in a public package
registry in early 2026. The end-to-end demo
(\texttt{examples/end\_to\_end/} in the repository) composes scanner,
prover, capability minting, gate, audit chain, and offline-verifiable
trust graph in a single script. The four pre-registered AgentDojo policy
files
(\texttt{paper/eval/policies/\{banking,slack,travel,workspace\}.json},
hash-anchored in \texttt{paper/eval/PRE-REGISTRATION.md}) authorised
\textbf{7,312} tool calls and denied \textbf{761} policy-violating calls
across the empirical evaluation reported in §6.2 --- an overall deny
rate of 9.4\% across 8,073 total gate decisions, consistent with the
AgentDojo benchmark's design in which most user × injection pairs
combine a legitimate multi-call user task with an adversarial
single-call attack attempt. We do not report a separate production
deployment in this paper; the integration cost is modest (the reference
adapter for AgentDojo is approximately 300 lines of glue), and we
anticipate that one or more public deployments will report their numbers
independently. The repository, Lean development, benchmark harness, and
policy files are intended as a complete reproducibility package; readers
can re-run any number in §6 against their own hardware and model
credentials.

\subsubsection{6.5 Negative Results}\label{negative-results}

\textbf{On the residual non-zero ASR in §6.2 and §6.2.2.} Every VCD row
across the LLM-driven evaluation with a non-zero ASR reports exactly one
attack success out of 144 banking pairs (0.7\% = 1/144). In every case
--- \texttt{vcd\_text}, \texttt{vcd\_full}, and \texttt{vcd\_cap\_only}
on all three models (v3.2, v4-flash, v4-pro), under all three attack
families (\texttt{important\_instructions}, \texttt{direct},
\texttt{ignore\_previous}) --- the single ``successful'' attack is the
same cell: \textbf{user\_task\_15 × injection\_task\_4}. This is the
known-legitimate IBAN coincidence flagged in §6.2.1: the user explicitly
asks the agent to schedule a rent payment to their landlord, the
landlord's IBAN happens to be the same IBAN the benchmark's adversarial
injection\_task\_4 designates as the attacker's target, and AgentDojo's
security oracle marks the cell \texttt{security=True} because \emph{a}
transfer to that IBAN occurred --- without distinguishing between the
user-authorised transfer (which VCD correctly admits, with a valid
capability token issued from the user's prompt) and an attacker-induced
one. The gate did not fail; the benchmark oracle cannot disambiguate.

We deliberately do not exclude this cell from the LLM-driven numbers.
The honest measurement is what the benchmark reports against the policy
we pre-registered; the §6.2.1 static analysis carries the structural
claim that the gate denies every attacker-controlled cell. In aggregate
across the \textbf{720} banking pairs measured under VCD full-stack and
its ablations (3 models × 1 attack + 1 model × 2 additional attacks,
each at 144 pairs), the gate admits every user-authorised call and
denies every attacker-induced call, with the single residual being the
benchmark's IBAN-collision artefact. Outside that one benchmark cell,
the LLM-driven gate-block rate on attacker-controlled tool calls is
\textbf{100\%}, matching the §6.2.1 static bound.

We separately report two classes where VCD \emph{legitimately} does not
eliminate attack success --- i.e., where the gate's design boundary is
the limit, not the oracle:

\begin{itemize}
\tightlist
\item
  \textbf{Tool-output exfiltration.} Attacks whose payload is the
  \emph{content} the agent eventually emits in its response --- e.g.,
  coercing the agent to print a secret retrieved from a legitimate tool
  call --- are not affected by the gate. The legitimate tool call was
  authorised by a valid token; the leak occurs in the model's text
  output. VCD makes no claim against this class. Composition with
  output-side defences is the obvious next step.
\item
  \textbf{Permitted-parameter side channels.} A schema permitting
  \texttt{amount} to be any number in
  \texttt{{[}0,\ 1\textbackslash{},000\textbackslash{},000{]}} permits
  an attacker who can coerce the agent into making \emph{one} legitimate
  \texttt{transfer\_funds} call to encode up to roughly 20 bits of
  information in the chosen amount. We do not eliminate this channel; we
  discuss schema tightening in Section 8.
\end{itemize}

We consider these honest concessions essential to the paper. The
strength of the headline claim depends on its narrowness.
